% Section 3: Likelihood and Identifiability
\section{Likelihood and Identifiability}
\label{sec:likelihood}

The system CDF~\eqref{eq:system-cdf} is a product of terms
$(1 - e^{-\lambda_j t})$, each of which is itself a difference of
exponentials.  Expanding this product via inclusion-exclusion yields the
analytical machinery that underlies all subsequent results: closed-form
densities, tractable likelihoods, and an elementary identifiability proof.

\subsection{Inclusion-exclusion expansion}
\label{sec:ie-expansion}

\begin{lemma}[Inclusion-exclusion expansion]\label{lem:ie}
  For rates $\lambda_1, \ldots, \lambda_m > 0$ and $t > 0$,
  \begin{equation}\label{eq:ie-expansion}
    \prod_{j=1}^m \bigl(1 - e^{-\lambda_j t}\bigr)
    = \sum_{S \subseteq \{1,\ldots,m\}} (-1)^{|S|}\, e^{-\lambda(S)\, t},
  \end{equation}
  where $\lambda(S) = \sum_{j \in S} \lambda_j$, with the convention
  $\lambda(\varnothing) = 0$.
\end{lemma}

\begin{proof}
By induction on~$m$; see \Cref{sec:appendix-proofs}.
\end{proof}

\begin{remark}\label{rmk:ie-terms}
The expansion~\eqref{eq:ie-expansion} has $2^m$ terms.  Each term
$e^{-\lambda(S) t}$ is an exponential with rate equal to a subset sum of
the~$\lambda_j$.  When the rates are distinct and no two subset sums
coincide---a generic condition---all $2^m$ terms have distinct rates.
This structural richness is ultimately what makes parallel systems
identifiable (\Cref{thm:identifiability}).
\end{remark}

\subsection{Density decomposition}
\label{sec:density-decomp}

\begin{lemma}[Density decomposition]\label{lem:density-decomp}
  The system density decomposes as
  \begin{equation}\label{eq:density-decomp}
    f_T(t; \blam) = \sum_{j=1}^m w_j(t; \blam),
  \end{equation}
  where the \emph{component weight function}
  \begin{equation}\label{eq:wj-def}
    w_j(t; \blam) = \lambda_j e^{-\lambda_j t} \prod_{i \neq j}
    \bigl(1 - e^{-\lambda_i t}\bigr)
  \end{equation}
  satisfies $w_j(t; \blam) \geq 0$ for all $t > 0$.  Moreover,
  \begin{equation}\label{eq:posterior-J}
    \frac{w_j(t; \blam)}{f_T(t; \blam)} = P(J = j \mid T = t; \blam),
  \end{equation}
  the conditional probability that component~$j$ was the last to fail,
  given system failure at time~$t$.
\end{lemma}

\begin{proof}
By direct computation of $P(T \in dt,\; J = j)$ via independence;
see \Cref{sec:appendix-proofs}.
\end{proof}

\begin{remark}\label{rmk:series-contrast-density}
The decomposition~\eqref{eq:density-decomp} is the parallel analogue
of the series decomposition $f_T(t) = \sum_j h_j(t) S_T(t)$, where
$h_j$ is the cause-specific hazard rate and $S_T$ is the system survival
function.  In the series case, hazard rate additivity
$h_T(t) = \sum_j h_j(t)$ provides a clean factorization.  The parallel
case has no such additivity; instead, the density is a sum of products,
with each term coupling one component's density to the CDFs of all
others.
\end{remark}

\subsection{Closed-form integrals}
\label{sec:closed-form-integral}

The weight functions $w_j$ can be integrated in closed form---a property
that is essential for handling censored and interval-censored system
observations without numerical quadrature.

\begin{lemma}[Closed-form integral of $w_j$]\label{lem:integral-wj}
  For $0 \leq a < b \leq \infty$ and all $\lambda_j > 0$,
  \begin{equation}\label{eq:integral-wj}
    \int_a^b w_j(t; \blam)\, dt
    = \lambda_j \sum_{S \subseteq \{1,\ldots,m\} \setminus \{j\}}
      (-1)^{|S|} \cdot
      \frac{e^{-r_S a} - e^{-r_S b}}{r_S},
  \end{equation}
  where $r_S = \lambda_j + \lambda(S) > 0$.
  In particular, when $b = \infty$,
  \begin{equation}\label{eq:integral-wj-tail}
    \int_a^\infty w_j(t; \blam)\, dt
    = \lambda_j \sum_{S \subseteq \{1,\ldots,m\} \setminus \{j\}}
      (-1)^{|S|} \cdot \frac{e^{-r_S a}}{r_S},
  \end{equation}
  and when $a = 0$ and $b = \infty$,
  \begin{equation}\label{eq:prob-last}
    P(J = j; \blam) = \int_0^\infty w_j(t; \blam)\, dt
    = \lambda_j \sum_{S \subseteq \{1,\ldots,m\} \setminus \{j\}}
      \frac{(-1)^{|S|}}{r_S},
  \end{equation}
  the unconditional probability that component~$j$ is the last to fail.
\end{lemma}

\begin{proof}
Apply \Cref{lem:ie} to the $m-1$ rates $\{\lambda_i : i \neq j\}$
appearing in the product within $w_j$:
\[
  w_j(t; \blam)
  = \lambda_j e^{-\lambda_j t} \prod_{i \neq j}
    \bigl(1 - e^{-\lambda_i t}\bigr)
  = \lambda_j \sum_{S \subseteq \{1,\ldots,m\} \setminus \{j\}}
    (-1)^{|S|}\, e^{-({\lambda_j + \lambda(S)})\, t}
  = \lambda_j \sum_S (-1)^{|S|}\, e^{-r_S t}.
\]
Since $r_S = \lambda_j + \lambda(S) \geq \lambda_j > 0$, each exponential
term is integrable on $[a, b]$:
\[
  \int_a^b e^{-r_S t}\, dt = \frac{e^{-r_S a} - e^{-r_S b}}{r_S}.
\]
Linearity of integration gives~\eqref{eq:integral-wj}.  The special
cases~\eqref{eq:integral-wj-tail} and~\eqref{eq:prob-last} follow by
substituting $b = \infty$ (so $e^{-r_S b} = 0$) and then $a = 0$
(so $e^{-r_S a} = 1$).
\end{proof}

\begin{remark}\label{rmk:integral-tractability}
\Cref{lem:integral-wj} is the key computational enabler for censored
likelihoods.  For right-censored system observations ($T > \tau$), the
survival probability $1 - F_T(\tau) = \int_\tau^\infty f_T(t)\, dt =
\sum_j \int_\tau^\infty w_j(t)\, dt$ reduces to a finite sum of
exponentials via~\eqref{eq:integral-wj-tail}.  For interval-censored
observations ($a < T \leq b$), the probability
$F_T(b) - F_T(a) = \sum_j \int_a^b w_j(t)\, dt$ likewise has a closed
form via~\eqref{eq:integral-wj}.  No numerical quadrature is needed.
\end{remark}

\subsection{Log-likelihood}
\label{sec:log-likelihood}

Suppose we observe $n$ system lifetimes, some exact and some censored.
Let $\mathcal{E}$, $\mathcal{R}$, and $\mathcal{I}$ denote the index
sets for exact, right-censored, and interval-censored observations,
respectively.  The observed-data log-likelihood is
\begin{equation}\label{eq:loglik}
  \ell(\blam) = \sum_{i \in \mathcal{E}} \ln f_T(t_i; \blam)
  + \sum_{i \in \mathcal{R}} \ln \bigl[1 - F_T(\tau_i; \blam)\bigr]
  + \sum_{i \in \mathcal{I}} \ln \bigl[F_T(b_i; \blam) - F_T(a_i; \blam)\bigr],
\end{equation}
where each term has a closed-form expression via the inclusion-exclusion
expansion:
\begin{align}
  f_T(t; \blam) &= \sum_{j=1}^m w_j(t; \blam)
  = \sum_{j=1}^m \lambda_j \sum_{S \subseteq \{1,\ldots,m\} \setminus \{j\}}
    (-1)^{|S|}\, e^{-r_S t},
  \label{eq:fT-ie} \\
  F_T(t; \blam) &= \sum_{S \subseteq \{1,\ldots,m\}}
    (-1)^{|S|}\, e^{-\lambda(S)\, t},
  \label{eq:FT-ie} \\
  F_T(b) - F_T(a) &= \sum_{S \subseteq \{1,\ldots,m\}}
    (-1)^{|S|}\, \bigl(e^{-\lambda(S)\, a} - e^{-\lambda(S)\, b}\bigr).
  \label{eq:interval-ie}
\end{align}

\noindent
Each evaluation of $\ell(\blam)$ requires $O(n \cdot m \cdot 2^{m-1})$
arithmetic operations.  The exponential scaling in~$m$ limits direct
application to systems with a moderate number of components
(say $m \leq 15$--$20$), but this covers many practical redundant
systems.

When all $n$ observations are exact ($\mathcal{R} = \mathcal{I} = \varnothing$),
the log-likelihood reduces to
\begin{equation}\label{eq:loglik-exact}
  \ell(\blam) = \sum_{i=1}^n \ln\!\Biggl[\,
    \sum_{j=1}^m w_j(t_i; \blam) \Biggr],
\end{equation}
the Scheme~0 (black-box) log-likelihood.

\subsection{Identifiability}
\label{sec:identifiability}

The central theoretical result of this paper is that the distribution of
the parallel system lifetime uniquely determines all component rates.

\begin{theorem}[Identifiability of exponential parallel systems]
\label{thm:identifiability}
  Let $T_1, \ldots, T_m$ be independent with $T_j \sim \mathrm{Exp}(\lambda_j)$,
  $\lambda_j > 0$, and let $T = \max(T_1, \ldots, T_m)$.  The distribution
  of~$T$ uniquely determines the multiset $\{\lambda_1, \ldots, \lambda_m\}$.

  That is, if $\blam = (\lambda_1, \ldots, \lambda_m)$ and
  $\boldsymbol{\mu} = (\mu_1, \ldots, \mu_m)$ both generate the same
  distribution for $T = \max_j T_j$, then
  $\{\lambda_1, \ldots, \lambda_m\} = \{\mu_1, \ldots, \mu_m\}$ as
  multisets.
\end{theorem}

\begin{proof}
The proof proceeds by induction on~$m$, extracting component rates one
at a time from the tail behavior of the reversed hazard rate.

\emph{Base case} ($m = 1$).
The CDF $F_T(t) = 1 - e^{-\lambda_1 t}$ determines $\lambda_1$
uniquely.

\emph{Inductive step.}  Assume the result holds for all parallel systems
with fewer than~$m$ components.  We show that $F_T$ determines the
multiset of rates.

\textbf{Step~1: Extract the smallest rate from the tail.}
The reversed hazard rate of~$T$
\citep[the natural hazard analogue for parallel systems;
see][]{BlockSavitsSingh1998} is
\[
  r_T(t) = \frac{f_T(t)}{F_T(t)}
  = \frac{d}{dt} \ln F_T(t)
  = \sum_{j=1}^m \frac{\lambda_j e^{-\lambda_j t}}{1 - e^{-\lambda_j t}}.
\]
For large~$t$, each summand admits the geometric series expansion
\[
  \frac{\lambda_j e^{-\lambda_j t}}{1 - e^{-\lambda_j t}}
  = \lambda_j e^{-\lambda_j t} \sum_{k=0}^\infty e^{-k\lambda_j t}
  = \lambda_j e^{-\lambda_j t} + O\bigl(e^{-2\lambda_j t}\bigr).
\]
Let $\lambda_* = \min_j \lambda_j$ denote the smallest rate, with
multiplicity~$\kappa$ (the number of components having rate~$\lambda_*$).
As $t \to \infty$,
\begin{equation}\label{eq:rhr-asymp}
  r_T(t) = \kappa\, \lambda_*\, e^{-\lambda_* t}
  + O\bigl(e^{-(\lambda_* + \delta) t}\bigr),
\end{equation}
where $\delta = \min\{\lambda_j - \lambda_* : \lambda_j > \lambda_*\} > 0$
if $\kappa < m$, and the remainder is exponentially smaller than the
leading term.  Since $r_T(t)$ is determined by $F_T$, we extract
$\lambda_*$ as
\[
  \lambda_* = \lim_{t \to \infty} \frac{-\ln r_T(t)}{t},
\]
and the multiplicity as
\[
  \kappa = \frac{1}{\lambda_*} \lim_{t \to \infty} r_T(t)\, e^{\lambda_* t}.
\]
Both limits exist and are uniquely determined by~$F_T$.

\textbf{Step~2: Factor out the identified components.}
Define
\[
  G(t) = \frac{F_T(t)}{\bigl(1 - e^{-\lambda_* t}\bigr)^\kappa}.
\]
Since $F_T(t) = (1 - e^{-\lambda_* t})^\kappa \prod_{j:\, \lambda_j > \lambda_*}
(1 - e^{-\lambda_j t})$, we have
\[
  G(t) = \prod_{j:\, \lambda_j > \lambda_*} \bigl(1 - e^{-\lambda_j t}\bigr),
\]
which is the CDF of $\max\{T_j : \lambda_j > \lambda_*\}$---a parallel system
with $m - \kappa < m$ components.  Since $F_T$ and $\lambda_*$ are known, $G$
is known.

\textbf{Step~3: Apply the induction hypothesis.}
By the induction hypothesis, the multiset $\{\lambda_j : \lambda_j > \lambda_*\}$
is uniquely determined by~$G$.  Combined with the $\kappa$ copies of~$\lambda_*$
identified in Step~1, the full multiset $\{\lambda_1, \ldots, \lambda_m\}$ is
determined.
\end{proof}

\begin{remark}\label{rmk:identifiability-role-of-independence}
The proof uses independence in two places: to factor $F_T$ as a product and
to decompose the reversed hazard rate as a sum.  Under dependence (e.g.,
copula models), identifiability may fail; see \citet{BasuGhosh1978} for
conditions and counterexamples.
\end{remark}

\subsection{Identifiability duality}
\label{sec:id-duality}

\begin{corollary}[Identifiability duality]\label{cor:duality}
  For independent exponential components with rates
  $\lambda_1, \ldots, \lambda_m > 0$:
  \begin{enumerate}
    \item[\textup{(a)}] \textbf{Series} ($T = \min_j T_j$).  The distribution
      of~$T$ determines only the total rate
      $\Lambda = \sum_{j=1}^m \lambda_j$, not the individual rates.
    \item[\textup{(b)}] \textbf{Parallel} ($T = \max_j T_j$).  The distribution
      of~$T$ determines the individual rates
      $\{\lambda_1, \ldots, \lambda_m\}$ without any additional information.
  \end{enumerate}
\end{corollary}

\begin{proof}
(a) The series survival function is
$\bar{F}_T(t) = \prod_j e^{-\lambda_j t} = e^{-\Lambda t}$, which depends
on~$\blam$ only through~$\Lambda$.  Any partition of~$\Lambda$ into~$m$
positive parts yields an observationally equivalent model.  This is
the classical non-identifiability result of \citet{Tsiatis1975}.

(b) \Cref{thm:identifiability}.
\end{proof}

\begin{remark}\label{rmk:duality-information}
The duality extends to the information-theoretic level.  The Fisher
information from system data alone has rank~$1$ for series systems (all
information concentrates in the direction $(1, \ldots, 1)$ in parameter
space) and full rank for parallel systems (\Cref{thm:fisher}).  The
incomplete-data problems are structural duals: series systems need
cause-of-failure data to separate the~$\lambda_j$; parallel systems are
already identified, but left-censoring of fast components degrades
estimation efficiency.
\end{remark}

\subsection{Fisher information}
\label{sec:fisher}

\begin{theorem}[Positive definiteness of Fisher information]
\label{thm:fisher}
  Under Scheme~0, the Fisher information matrix for a single observation
  of $T = \max(T_1, \ldots, T_m)$ with $T_j \sim \mathrm{Exp}(\lambda_j)$,
  \begin{equation}\label{eq:fisher-def}
    I(\blam)_{jk} = \E\!\left[
      \frac{\partial \ln f_T(T; \blam)}{\partial \lambda_j}
      \cdot \frac{\partial \ln f_T(T; \blam)}{\partial \lambda_k}
    \right],
  \end{equation}
  is positive definite for all $\blam \in \Rset_{>0}^m$ with distinct
  rates.
\end{theorem}

\begin{proof}
The Fisher information matrix is the variance-covariance matrix of the
score vector $\nabla_{\blam} \ln f_T(T; \blam)$, hence positive
semidefinite.  We must show it is nonsingular.

Suppose $I(\blam) v = 0$ for some $v \in \Rset^m$.  Then
\[
  \Var\!\bigl(v^\top \nabla_{\blam} \ln f_T(T; \blam)\bigr)
  = v^\top I(\blam)\, v = 0,
\]
so $v^\top \nabla_{\blam} \ln f_T(T; \blam) = 0$ almost surely.  That
is,
\[
  \sum_{j=1}^m v_j \frac{\partial}{\partial \lambda_j}
  \ln f_T(t; \blam) = 0
  \qquad \text{for almost every } t > 0,
\]
where the almost-sure statement under~$P$ transfers to
Lebesgue-almost-every~$t$ because $f_T(t; \blam) > 0$ for all $t > 0$.
This implies that the model is locally
non-identifiable in the direction~$v$.

By \Cref{thm:identifiability}, the model
$\{f_T(\cdot; \blam) : \blam \in \Rset_{>0}^m\}$ is globally
identifiable up to permutation of the rates.  For distinct rates, the
permutation equivalence class $\{\blam_\sigma : \sigma \in S_m\}$ is a
finite (hence discrete) set, through which no smooth curve can pass.
Therefore there are no
nontrivial directions of local non-identifiability, and $v = 0$.
\end{proof}

\begin{remark}\label{rmk:fisher-explicit}
While the positive definiteness of $I(\blam)$ is guaranteed by
\Cref{thm:fisher}, closed-form expressions for its entries require
integration of products of score components against $f_T$---integrals
that do not simplify to elementary form.  In practice, $I(\blam)$ is
computed numerically, either by Monte Carlo integration or by the
observed information matrix $-\nabla^2 \ell(\hat{\blam})$ evaluated
at the MLE.

The partial derivatives needed for the score are available analytically
via the inclusion-exclusion expansion.  From~\eqref{eq:wj-def},
differentiating $w_j$ with respect to $\lambda_j$ and $w_k$ ($k \neq j$)
with respect to $\lambda_j$ gives:
\begin{align}
  \frac{\partial w_j(t)}{\partial \lambda_j}
  &= \sum_{S \subseteq \{1,\ldots,m\} \setminus \{j\}}
    (-1)^{|S|}\, \bigl(1 - \lambda_j t\bigr)\, e^{-r_S t},
  \label{eq:dwj-dlamj} \\[6pt]
  \frac{\partial w_k(t)}{\partial \lambda_j}
  &= -\lambda_k \sum_{\substack{S \subseteq \{1,\ldots,m\} \setminus \{k\} \\
    j \in S}} (-1)^{|S|}\, t\, e^{-({\lambda_k + \lambda(S)})\, t},
  \qquad k \neq j.
  \label{eq:dwk-dlamj}
\end{align}
These are finite sums of exponentials, so the score and Hessian of
$\ell(\blam)$ are analytically tractable even though the expected Fisher
information is not.
\end{remark}
